\lyseis{Β1.10}{Απόσταση σημείου από ευθεία και απόσταση παραλλήλων}
\sol{1} Το κάθετο τμήμα $\mathrm{ΑΜ}$ είναι μικρότερο από κάθε πλάγιο.
\sol{2} Η απόσταση είναι το μήκος του κάθετου τμήματος $\mathrm{ΑΜ}$. Για σημείο της ευθείας είναι $0$.
\sol{3} Βγαίνει η ίδια απόσταση σε όλα τα σημεία.
\sol{4} $2{,}4+2{,}4=4{,}8$ cm.
\sol{5} Δύο παράλληλες προς την $(\varepsilon)$ σε απόσταση $2$ cm, μία σε κάθε ημιεπίπεδο.
\sol{6} α) Σωστό. β) Λάθος. γ) Σωστό. δ) Λάθος.
\sol{7} Όχι. Το λοξό μήκος ($3{,}6$ μ.) είναι μεγαλύτερο από το κάθετο ($3$ μ.), όπως πρέπει.
\sol{8} $60+60=120$ cm.
\sol{9} Όχι $3$ cm (το πλάγιο δεν μπορεί να είναι μικρότερο από το κάθετο $4$). $4$ cm μόνο αν το $\mathrm{Β}$ είναι το ίχνος της καθέτου. $6$ cm ναι.
\sol{10} $1{,}2$ μ. σε κάθε σημείο.
\sol{11} Ανάμεσα στις μεγάλες πλευρές $25$ μ. και ανάμεσα στις μικρές $40$ μ.
\sol{12} Για να φτάσεις από ένα σημείο στην ευθεία, κάθε πλάγιο δρόμος είναι «λοξότερος», άρα μεγαλύτερος από τον κάθετο, όπως ένα τεντωμένο νήμα που στρέφεται και μακραίνει.
\sol{13} Όταν οι δύο ευθείες ταυτίζονται.
\sol{14} Δύο παράλληλες ευθείες με $3+5=8$ cm: τα σημεία είναι σε μία ευθεία παράλληλη που απέχει $3$ cm από την πρώτη και $5$ από τη δεύτερη.
\sol{15} Μία: η μεσαία παράλληλη.
\sol{16} Ελεύθερη απάντηση.
\sol{17} Ελεύθερη απάντηση: μετράς κάθετα στον τοίχο.
